Tuning particle accelerators is an expensive and time consuming challenge that causes limitations in the findings of advanced research facilities. Bayesian Optimization (BO) has shown potential to be a promising approach for autonomous tuning of accelerators, yet its effectiveness diminishes for higher dimensional problems. Recent studies have proposed replacing the zero mean prior in BO with physics informed prior models to leverage physics domain knowledge. The Cheetah framework, a high-speed differentiable simulator, was demonstrated as a physics informed BO prior on a simple two dimensional FODO lattice. This project work extends that approach to a real world based five dimensional problem now: transverse beam parameter optimization at ARES Experimental Ares, a five dimensional problem involving three quadrupole magnets and two corrector magnets. This project implements a framework that incorporates Cheetah as the Gaussian Process mean function into the Xopt optimization toolset with trainable misalignment parameters that enable the prior to adjust when the simulator does not accurately represent the real system. Four optimization approaches are evaluated. Results from 10 trails show that physics informed BO performed significantly better than baseline techniques. Both Cheetah prior variants achieve 98.4\% improvement in beam error with mean absolute error of 0.006 mm, compared to standard BO with 88.9\% and 0.046 mm, and Nelder-Mead with 83.7\% and 0.067 mm. While standard BO only reaches 20\% and Nelder-Mead 0\%, both Cheetah variants managed to achieve a 100\% success rate in achieving the desired threshold (40 $\mu$m). Moreover, the matched prior converged the fastest, requiring only 4 steps (median) compared to 14 for mismatched prior.  These results validate the scalability of differentiable simulator priors simple and straightforward to realistic and complex accelerator tuning problems. 
